\scanpage{111}
\mspage{98}
% archive sheet number: 110

6. \uline{Application à la construction des dérivateurs}
(associés aux triples de Quillen spéciaux).
% slanted
\margin{C'est le résultat \uline{principal} \unsure{de|dans} \unsure{ma|la} \unsure{partie} $\overline{\underline{\mathrm{XVII}}}$ des notes.}

\uline{Théorème} 6.1. Soit \deleted{\unsure{$M$}} $(W,C,F)$ un triple
de Quillen clos spécial \deleted{\unsure{dans}} \inserted{sur la} catégorie $M$
(cf déf.~5.6. p.~110, et cor.~5.7). Pour
tout $I$ dans Cat, considérons la catégorie
\[
 M(I)=\underline{\mathrm{Hom}}(I^{\unsure{o|\circ}},M),\tag*{(6.1.1)}
\]
et \deleted{posons} considérons
\[
 W(I),\deleted{\unsure{C}}F(I)\subset\mathrm{Fl}\,M(I)\tag*{\deleted{(6.1.2)}}
\]
définis par
\[
\begin{aligned}
 &u:X\mathrel{\unsure{-|\longrightarrow}}Y\ \text{dans }M(I)\text{ est dans }\unsure{W(I)|W(F)}\\
 &\text{(resp. }F(I))\ \text{ssi }\ \forall i\in\mathrm{Ob}\,I,\quad u(i):X(i)\longrightarrow Y(i)\\
 &\text{est dans }W\ \text{(resp. }F).
\end{aligned}\tag*{(6.1.2)}
\]
\deleted{Alors} \inserted{a)} Il existe une structure de
Quillen close (néc. unique, \inserted{\unsure{avec} $C(I)=$}) dans $M(I)$,
ayant $W(I)$ et $F(I)$ comme \deleted{équiv. fai} \inserted{\unsure{ensembles de|ens. des}}
quasi-isomorphismes et de fibrations
respectivement. On a donc aussi
\[
\left\{\begin{aligned}
 TF(I)&=F(I)\cap W(I)=\{u\in\mathrm{Fl}(M(I))\mid\\
 &\hspace{1.25em}u(i)\in TF\quad\forall i\in\mathrm{Ob}\,I\},\unsure{.}\\
 C(I)&=TF(I)^{*}.
\end{aligned}\right.\tag*{(6.1.3)}
\]
Cette structure de Quillen est \uline{spéciale}.

